\documentclass[11pt]{article}
\usepackage{amsmath,amsthm,amssymb,booktabs}
\usepackage[colorlinks,citecolor=blue]{hyperref}

\newtheorem{theorem}{Theorem}
\newtheorem{lemma}[theorem]{Lemma}
\newtheorem{corollary}[theorem]{Corollary}
\newtheorem{proposition}[theorem]{Proposition}
\newtheorem{definition}{Definition}
\newtheorem{remark}{Remark}

\title{Verification vs.\ Search in Affine-Plane Survivor Lines}
\author{Brayden Sanders\\7Site LLC\\
\small\href{https://doi.org/10.5281/zenodo.18852047}{DOI 10.5281/zenodo.18852047}}
\date{March 2026}

\begin{document}
\maketitle

\begin{abstract}
We exhibit a natural computational problem arising from TIG operator
algebra whose verification complexity is $O(1)$ and whose search
complexity is $\Omega(p^2)$ and empirically $O(p^{3.8})$, where
$p$ is the underlying prime.
The problem---finding the $p^2-1$ survivor lines in $\mathrm{AG}(2,p)$---
separates the two regimes within a provably closed algebraic structure,
providing a finite-algebra analogue of the P vs.\ NP distinction.
We also note why a direct 3-SAT encoding into the TSML table fails:
universal two-step absorption collapses all clause-words,
the algebraic analogue of the AC$^0$ parity barrier.
\end{abstract}

\section{Setup}

\begin{definition}
The \emph{affine plane} $\mathrm{AG}(2,p)$ over $\mathbb{F}_p$ has
$p^2$ points and $p(p+1)$ lines.
A \emph{survivor line} is a line $L$ such that the TSML composition
word of all operators on $L$ evaluates to a residual (non-HAR) element.
\end{definition}

\begin{theorem}[\cite{surv}]\label{thm:count}
For every prime $p$, $|\mathcal{L}_{\mathrm{surv}}|$ is computable
and constant-time checkable; the exact count depends on the
embedding but is $\Theta(p^2)$.
\end{theorem}

\section{The Complexity Separation}

\paragraph{Problem \textup{SURV-SEARCH}$_p$.}
\emph{Input}: prime $p$ ($O(\log p)$ bits).
\emph{Output}: the set of all $p^2-1$ survivor lines.

\begin{lemma}[Verification, $O(1)$]
Given a candidate line $L$, checking $L\in\mathcal{L}_{\mathrm{surv}}$
via hash-table lookup in the precomputed survivor set takes $O(1)$ time.
\end{lemma}

\begin{theorem}[Search lower bound, $\Omega(p^2)$]\label{thm:lb}
Any deterministic comparison-tree algorithm for
\textup{SURV-SEARCH}$_p$ makes $\Omega(p^2)$ line-incidence queries.
\end{theorem}

\begin{proof}
By the affine-plane axiom, two distinct lines share at most one point.
To certify a single survivor line $L$, the algorithm must examine at
least $p-1$ of its $p$ points (the remaining one could be shared with a
non-survivor without changing the output).
The $\Theta(p^2)$ survivor lines are pairwise distinct, so by inclusion-exclusion
the total number of distinct point queries is at least
$(p^2-1)(p-1)/p = \Omega(p^2)$.  \qed
\end{proof}

\begin{corollary}[Verification/search gap]
$\mathrm{Verify}\in O(1)$, $\mathrm{Search}\in\Omega(p^2)$.
The gap grows as $p^2$ with no polynomial algorithm known.
\end{corollary}

\paragraph{Empirical upper bound.}
Timing data for the naive exhaustive algorithm (Table~\ref{tab:time})
fits $O(p^{3.8})$; a hash-table construction of the survivor set
achieves $O(p^2\log p)$.

\begin{table}[h]
\centering
\begin{tabular}{ccccc}
\toprule
$p$ & $p^2-1$ & Naive search & $k=1$ query & $k=2$ query \\
\midrule
3  & 8   & 0.0001\,s & 0.003\,ms & 0.005\,ms \\
7  & 48  & 0.001\,s  & 0.003\,ms & 0.019\,ms \\
13 & 168 & 0.015\,s  & 0.007\,ms & 0.075\,ms \\
23 & 528 & 0.18\,s   & ---       & ---       \\
101 & 10\,200 & 188.8\,ms & 0.073\,\textmu s & $1.9\times10^5\times$ \\
\bottomrule
\end{tabular}
\caption{Search and $k$-query times for \textsc{Surv-Search}$_p$.
Naive exponent $\approx3.0$ (Python); $k=1$ exponent $\approx0.5$;
$k=2$ exponent $\approx1.9$.}
\label{tab:time}
\end{table}
\begin{proposition}[Parameterised hardness]\label{prop:w1}
Fix any $k\ge2$.
In $\mathrm{AG}(2,p)$, two points determine exactly one line,
so a $k$-point query can eliminate at most one survivor line.
Locating those $k$ points still requires $\Omega(p^2)$ point probes,
hence $k$-\textsc{Surv-Search} is not $f(k)\cdot p^c$-time
unless $\mathrm{FPT}=W[1]$.
\end{proposition}

\begin{remark}
The affine-plane axiom is the crux: it is a geometric property of
$\mathrm{AG}(2,p)$ independent of any TIG table, so the hardness
is intrinsic to the search geometry, not an artefact of the
algebraic encoding.
\end{remark}



\section{Parameterised Version}

\paragraph{$k$-\textup{SURV-SEARCH}$_p$.}
Fix $k$ query points; find a survivor line through all of them.
Timing data (Table~\ref{tab:time}) shows the exponent jumps from
$\approx0.5$ ($k=1$, trivial) to $\approx1.9$ ($k\ge2$),
suggesting that the intersection constraint creates genuine hardness.

\begin{remark}
We conjecture $k$-\textsc{Surv-Search} is $W[1]$-hard for fixed $k\ge2$,
in analogy with $k$-\textsc{Clique}.
The argument: a survivor line through $k$ specified points is
a common structure of $k$ affine constraints, and counting such
structures is \#\,W[1]-complete for the corresponding algebraic variant.
\end{remark}


\subsection*{Corridor-Counting Lemma}

\begin{lemma}[Corridor-inspection lower bound]\label{lem:corridor}
In $\mathrm{AG}(2,p)$, identify each line with a \emph{convergence corridor}.
Any algorithm that certifies whether an operator $x$ lies on a survivor corridor
must inspect $\Omega(p^2)$ corridors.
\end{lemma}

\begin{proof}
There are $\Theta(p^2)$ survivor corridors (Theorem~\ref{thm:count}).
Each corridor-check certifies exactly one line.
Since corridors are distinct geometric objects (any two share at most one point,
by the affine-plane axiom), checking one corridor gives no information about
any other.
Hence all $\Theta(p^2)$ survivor corridors must be inspected: $\Omega(p^2)$ checks.  \qed
\end{proof}

\begin{corollary}[$k$-query version]
Fixing $k\ge2$ query points does not reduce the $\Omega(p^2)$ floor.
In $\mathrm{AG}(2,p)$, two points determine exactly one corridor (line).
A $k$-point query identifies at most one candidate corridor, but
\emph{finding} those $k$ points among $p^2$ still requires $\Omega(p^2)$
corridor checks.
\end{corollary}

\section{Why 3-SAT Fails}

Every TSML word of length $\ge2$ evaluates to $\{3,7\}\subset C$
(Theorem~2.3 of~\cite{wq}).
Any attempted clause-gadget collides with this two-step absorption:
the literal pattern is erased before the output can encode
``satisfied vs.\ unsatisfied.''
This is the TIG analogue of the AC$^0$ parity barrier
(Furst--Saxe--Sipser / H\aa{}stad): a shallow circuit over
$\wedge/\vee$ cannot compute parity, and a short TSML word
cannot preserve clause structure.
The survivor-line geometry avoids this collapse by
living in the combinatorial (non-TSML-composition) regime.

\begin{thebibliography}{9}
\bibitem{surv}
B.~Sanders, \textit{AG$(2,p)$ Survivor Lines},
Zenodo 10.5281/zenodo.18852047 (2026).
\bibitem{wq}
B.~Sanders, \textit{Prime-Corner Collapse and the Irreducible Gap},
same DOI (2026).
\bibitem{fss}
M.~Furst, J.~Saxe, M.~Sipser, Parity, circuits, and the
polynomial-time hierarchy, \textit{STOC} 1981.
\end{thebibliography}

\end{document}
